% ===================================================================
%  तालिका सं. 6 — घर के भीतर। — साज-सामान। — भोजन। — रोशनी।
%  Traduction 2026 d'apres texto/io, controlee sur texto/fr et
%  texto/en. Consigne : voir texto/hi/10-talika-01.tex.
%
%  Deux notes, toutes deux marquees « (*) », et deux appels : celui de
%  « babo » au bloc c1-12-1, celui de « lambrequino » au bloc c2-01-1.
%  L'ordre les departage, comme dans les autres colonnes.
%
%  LA FEMME DE CHAMBRE N'A PAS DE GROS PLAN, et c'est normal. L'ido du
%  bloc c1-06-1 porte « la chambristino (6) », le francais « la femme
%  de chambre (16) » : c'est l'ido qui se trompe, son propre numero 6
%  etant le savon du bloc c1-02-1. Le hindi suit donc le francais et
%  ecrit (16) ; la planche ne porte pas de 16 releve, et le renvoi
%  reste du texte ordinaire. C'est un des trois ecarts declares dans
%  outils/renvoji.py.
% ===================================================================

%%K t06-apar-1 apar
\VUsaut{6.0mm}
\VUtitre{15.0pt}{60}{तालिका सं. 6}
\VUsaut{1.4mm}
\VUtitre{11.4pt}{0}{घर के भीतर, साज-सामान, भोजन,}
\VUsaut{0.4mm}
\VUtitre{11.4pt}{0}{रोशनी।}
\VUsaut{2.2mm}

%%K t06-apar-2 apar
घर बनकर तैयार है। जॉन के चाचा अपने नए घर में आ बसे हैं, जिसकी बनावट हम जानते ही
हैं। अब देखें कि उन्होंने उसे कैसे सजाया है। पहले हम देखेंगे :

%%K t06-tit-1 sub
\VUpk{10.2pt}{40}{शयन-कक्ष।}
\VUsaut{3.0mm}

%%K t06-c1-01-1 p
1. --- हम कुछ बेमौके आ गए हैं, क्योंकि नौकरानी \VUgras{कमरे} की सफ़ाई में लगी
है।

%%K t06-c1-02-1 p
2. --- \VUgras{हाथ-मुँह धोने की चौकी} \textsuperscript{(1)} पर इधर-उधर वे चीज़ें
पड़ी हैं जिनसे नहाया, कंघी की और तैयार हुआ जाता है। वे हैं \VUgras{चिलमची}
\textsuperscript{(2)} और \VUgras{पानी का जग} \textsuperscript{(3)},
\VUgras{बाल का ब्रश} \textsuperscript{(4)}, \VUgras{कंघी} \textsuperscript{(5)},
\VUgras{साबुन} \textsuperscript{(6)} और \VUgras{स्पंज} \textsuperscript{(7)}, और
अंत में कुल्ला करने के पानी की एक छोटी \VUgras{शीशी} \textsuperscript{(8)}।

%%K t06-c1-03-1 p
3. --- ज़मीन पर, चौकी के सामने, यह रही गंदे पानी की \VUgras{बालटी}
\textsuperscript{(9)}।

%%K t06-c1-04-1 p
4. --- \VUgras{बाहरी कमरे} \textsuperscript{(10)} में हमें कुछ \VUgras{खूँटियाँ}
\textsuperscript{(13)} दिखाई देती हैं और दो \VUgras{छाते} \textsuperscript{(11)},
जो \VUgras{छाता-दान} \textsuperscript{(12)} में हैं।

%%K t06-c1-05-1 p
5. --- दरवाज़े की दूसरी ओर \VUgras{शीशेदार अलमारी} \textsuperscript{(14)} खड़ी है,
जिसमें \VUgras{कपड़े} \textsuperscript{(15)} रखे रहते हैं।

%%K t06-c1-06-1 p
6. --- \VUgras{नौकरानी} \textsuperscript{(16)} \VUgras{पलंग} \textsuperscript{(17)}
लगा रही है, इसी बहाने उसके अलग-अलग भाग देख लें। \VUgras{चारपाई}
\textsuperscript{(20)} पर एक अच्छा \VUgras{कमानीदार गद्दा} \textsuperscript{(21)} है,
जिस पर जस्तीन अभी ऊन का \VUgras{गद्दा} \textsuperscript{(22)} बिछाएगी। फिर वह
कुरसियों से दो \VUgras{चादरें} \textsuperscript{(23)} और \VUgras{कंबल}
\textsuperscript{(24)} उठाएगी। फिर वह पलंग के सिरहाने \VUgras{लंबा तकिया}
\textsuperscript{(25)} और \VUgras{तकिया} \textsuperscript{(26)} रखेगी, जो अभी
\VUgras{रजाई} \textsuperscript{(27)} के साथ \VUgras{पलंग के आगे की दरी}
\textsuperscript{(32)} पर पड़े हैं। वह \VUgras{परदों} \textsuperscript{(19)} और
\VUgras{छत्र} \textsuperscript{(18)} पर पंख का झाड़न फेरती है, और आधा काम हो जाता
है।

%%K t06-c1-07-1 p
7. --- फिर उसे \VUgras{पलंग के पास की मेज़} \textsuperscript{(28)} थोड़ी सँवारनी
होगी, \VUgras{अलार्म घड़ी} \textsuperscript{(29)} में चाबी भरनी होगी,
\VUgras{रात की बत्ती} \textsuperscript{(30)} तैयार करनी होगी, और
\VUgras{मोमबत्ती} \textsuperscript{(31)} हटाकर शमादान साफ़ करना होगा।

%%K t06-tit-2 sub
\VUsaut{5.0mm}
\VUpk{10.2pt}{40}{सिलाई की टोकरी और अँगीठी का सामान।}
\VUsaut{3.0mm}

%%K t06-c1-08-1 p
8. --- \VUgras{सिलाई की टोकरी} \textsuperscript{(34)}, जो \VUgras{तिपाई}
\textsuperscript{(33)} के पास रखी है, उसे भी सँवारने की बड़ी ज़रूरत है। उसे \VUgras{कढ़ाई}
\textsuperscript{(35)} तह करनी होगी, \VUgras{सिलाई की मेज़} \textsuperscript{(38)} में
\VUgras{अँगुश्ताना}, \VUgras{धागा} \textsuperscript{(36)}, \VUgras{सुइयाँ}
\textsuperscript{(37)} और \VUgras{कैंची} \textsuperscript{(39)} रखनी होंगी, और मेज़
का ढक्कन \VUgras{चाबी} \textsuperscript{(40)} से बंद करना होगा।

%%K t06-c1-09-1 p
9. --- बीच में रखी \VUgras{तिपाई-मेज़} \textsuperscript{(41)} पर जस्तीन
\VUgras{सुराही} \textsuperscript{(42)} का पानी बदलेगी। वह \VUgras{चीनीदान}
\textsuperscript{(43)} में \VUgras{चीनी} डालेगी, \VUgras{चीनी की चिमटी}
\textsuperscript{(44)} साफ़ करेगी और \VUgras{कपड़े का ब्रश} \textsuperscript{(46)}
हटा देगी।

%%K t06-c1-10-1 p
10. --- कमरे का दाहिना भाग उसे कम काम देगा, क्योंकि \VUgras{दिन का पलंग}
\textsuperscript{(47)} और उसका \VUgras{गद्दी} \textsuperscript{(48)} अपनी जगह पर हैं,
और चूँकि मौसम ठंडा नहीं है, \VUgras{अँगीठी} \textsuperscript{(49)} के सामान में
कुछ हिलाया नहीं गया। \VUgras{बेलचा} \textsuperscript{(50)}, \VUgras{चिमटा}
\textsuperscript{(51)}, \VUgras{लकड़ी के टेक} \textsuperscript{(55)} और
\VUgras{राख का परदा} \textsuperscript{(56)} साफ़ हैं; \VUgras{आग का परदा}
\textsuperscript{(54)} हिलाया नहीं गया और \VUgras{लकड़ी का संदूक}
\textsuperscript{(52)} \VUgras{जलावन} \textsuperscript{(53)} से भरा है।

%%K t06-noto-1 noto
\VUnotes{143.2mm}{%
(*) Babo = नन्हा बच्चा; अंग्रेज़ी \textit{baby}, फ़्रेंच \textit{bébé},
इतालवी \textit{bambino}।%
}

%%K t06-noto-2 noto
\VUnotes{143.2mm}{%
(*) Lambrequino = फ़्रेंच \textit{lambrequin}, स्पेनी
\textit{lambrequines}, पुर्तगाली \textit{lambrequins}, और जर्मन तथा अंग्रेज़ी
में भी \textit{lambrequins} --- अर्थात् जर्मन, अंग्रेज़ी, फ़्रेंच, पुर्तगाली
और स्पेनी, सब में एक ही।%
}

%%K t06-c1-11-1 p
11. --- उसे अभी \VUgras{अँगीठी की घड़ी} \textsuperscript{(57)},
\VUgras{शमादान} \textsuperscript{(58)} और \VUgras{सजावट की तश्तरियाँ}
\textsuperscript{(59)} झाड़नी होंगी, फिर \VUgras{आईना} \textsuperscript{(60)} पोंछना
होगा।

%%K t06-c1-12-1 p
12. --- रही बच्चे का \VUgras{पालना} \textsuperscript{(61)} (वह babo (*)), वह तो
पहले से तैयार है।

%%K t06-tit-3 sub
\VUsaut{5.0mm}
\VUpk{10.2pt}{40}{बैठक।}
\VUsaut{3.0mm}

%%K t06-c2-01-1 p
1. --- अब यह रही बैठक। यह बहुत सुंदर \VUgras{कमरा} है। दाहिने कोने की अँगीठी एक
नाज़ुक \VUgras{मूर्ति} \textsuperscript{(1)} और दो सुंदर \VUgras{फूलदानों}
\textsuperscript{(3)} से सजी है, और मख़मल की \VUgras{झालर} \textsuperscript{(2)} (*)
से ढकी है।

%%K t06-c2-02-1 p
2. --- अँगीठी की चिनगारियाँ कालीन में आग न लगा दें, इसलिए उसके सामने ताँबे का
एक \VUgras{आग-जँगला} \textsuperscript{(4)} रखा गया है।

%%K t06-c2-03-1 p
3. --- उसके पास एक सजीली \VUgras{लेखन-मेज़} \textsuperscript{(5)} खुली खड़ी है।
यहीं \VUgras{घर की मालकिन} \textsuperscript{(23)} अपने पत्र लिखती हैं।

%%K t06-c2-04-1 p
4. --- खिड़की पर \VUgras{जालीदार परदे} \textsuperscript{(6)} लगे हैं, जिन्हें \VUgras{डोरियाँ}
\textsuperscript{(8)} बाँधे रखती हैं, और वे डोरियाँ \VUgras{खूँटियों}
\textsuperscript{(7)} में अटकी हैं, और उन पर भारी कपड़े के परदे, जिनके ऊपर एक
\VUgras{झालर-पट्टी} \textsuperscript{(9)} है।

%%K t06-c2-05-1 p
5. --- खिड़की के ताक में एक \VUgras{गमला-चौकी} \textsuperscript{(11)} में हरा
\VUgras{पौधा} \textsuperscript{(10)} है, एक शानदार ताड़, जिसके पत्ते लगभग
\VUgras{तेल के लैंप} \textsuperscript{(12)} तक पहुँचते हैं।

%%K t06-c2-06-1 p
6. --- \VUgras{लैंप-छाया} \textsuperscript{(13)} मारिया ने बनाई है, वही मोहक
\VUgras{युवती} \textsuperscript{(14)} जो \VUgras{पियानो} \textsuperscript{(15)} पर
बैठी है। \VUgras{चौकी} \textsuperscript{(16)} पर सीधी बैठी वह एक धुन का अभ्यास कर
रही है। वह बहुत कुशल और बहुत सलीकेदार है, जैसा उसकी
\VUgras{स्वरलिपि की अलमारी} \textsuperscript{(17)} से पता चलता है : उसमें पूरा
क्रम है।

%%K t06-tit-4 sub
\VUsaut{5.0mm}
\VUpk{10.2pt}{40}{शरारती।}
\VUsaut{3.0mm}

%%K t06-c2-07-1 p
7. --- उसका भाई पॉल एक \VUgras{किताब} \textsuperscript{(19)} ढूँढ़ रहा है,
\VUgras{किताबों की अलमारी} \textsuperscript{(18)} में। वह ऐसा \VUgras{लड़का}
\textsuperscript{(20)} है जो चंचल और बिलकुल बेकाबू है। बहुत ऊँची रखी किताब तक
पहुँचने के लिए अलमारी से लटकते हुए वह उसके ऊपर रखी \VUgras{प्रतिमा}
\textsuperscript{(21)} गिरा देने का ख़तरा उठा रहा है।

%%K t06-c2-08-1 p
8. --- यह शरारत करने के लिए वह इस बात का लाभ उठा रहा है कि उसकी माँ एक सहेली से
बातों में लगी हैं, एक \VUgras{महिला} \textsuperscript{(22)} जो कुछ दिन उनके साथ
बिताने आई हैं। माँ \VUgras{सोफ़े} \textsuperscript{(24)} पर बैठी हैं,
\VUgras{आरामकुरसी} \textsuperscript{(25)} के पास, और उनकी सहेली उस
\VUgras{कुरसी} \textsuperscript{(27)} पर, जिसकी हमें केवल \VUgras{पीठ}
\textsuperscript{(26)} दिखाई देती है।

%%K t06-c2-09-1 p
9. --- दीवार पर टँगे बड़े \VUgras{फ़्रेम} \textsuperscript{(30)} में एक रिश्तेदार
का \VUgras{चित्र} \textsuperscript{(29)} है। मेज़ पर पड़ी \VUgras{एल्बम}
\textsuperscript{(32)} में बाकी परिवार की \VUgras{तस्वीरें} \textsuperscript{(33)}
इकट्ठी हैं। बच्चे अभी थोड़ी देर पहले उसे उलट-पलट रहे थे, क्योंकि
\VUgras{कुरसियाँ} \textsuperscript{(31)} अब भी मेज़ के चारों ओर जमा हैं।

%%K t06-c2-10-1 p
10. --- चीनी मिट्टी का \VUgras{चीनी फूलदान} \textsuperscript{(34)}, जो
\VUgras{काग़ज़ काटने की छुरी} \textsuperscript{(35)} के पास रखा है, मारिया को उसके नाम-दिवस पर मिला था,
और कमरे का कोना भरता \VUgras{तह होने वाला परदा} \textsuperscript{(36)} उसके चाचा
चीन से लाए थे।

%%K t06-c2-11-1 p
11. --- सुंदर \VUgras{चित्रों} \textsuperscript{(37)} से जगमगाती, जो
\VUgras{दीवार} \textsuperscript{(42)} पर टँगे हैं, छत के \VUgras{झाड़-फ़ानूस} \textsuperscript{(38)}
से रोशन, जिसके \VUgras{बिजली के लैंप} \textsuperscript{(39)} शाम को इतने प्रसन्न
चमकते हैं, यह बैठक बड़ी सुहावनी लगती है। फ़र्श पर बिछा कोमल \VUgras{कालीन}
\textsuperscript{(40)} और इतनी काम की \VUgras{बिजली की घंटी} \textsuperscript{(41)}
उसे पूरी तरह आरामदेह बना देते हैं।

%%K t06-tit-5 sub
\VUsaut{3.6mm}
\VUpk{10.2pt}{40}{भोजन-कक्ष।}
\VUsaut{3.0mm}

%%K t06-c3-01-1 p
1. --- भोजन की घंटी बज चुकी है। मारिया को छोड़ सारा परिवार मेज़ के चारों ओर
जमा है। भोजन-कक्ष एक बढ़िया मिट्टी के \VUgras{चूल्हे} \textsuperscript{(1)} से
सुखद गरम है, जिसकी \VUgras{नली} \textsuperscript{(2)} पतली है।

%%K t06-c3-02-1 p
2. --- इस कमरे में सामान अधिक नहीं। दीवार के साथ, \VUgras{तिपाई}
\textsuperscript{(4)} के ऊपर, एक सजावटी छोटा \VUgras{फ़व्वारा} \textsuperscript{(3)}
टँगा है। पास ही \VUgras{सजावटी अलमारी} \textsuperscript{(5)} है, जिस पर ठंडे
व्यंजन, मिठाइयाँ और वे बरतन रखे जाते हैं जिनकी इस समय ज़रूरत नहीं।

%%K t06-c3-03-1 p
3. --- सो ऊपर के ख़ाने में हमें एक \VUgras{भुना मुर्ग़ा} \textsuperscript{(7)}
दिखाई देता है, एक \VUgras{मछली} \textsuperscript{(8)} और एक \VUgras{कटोरा}
\textsuperscript{(10)} \VUgras{फलों} \textsuperscript{(9)} का। नीचे
\VUgras{पनीर} \textsuperscript{(11)}, \VUgras{रोटी} \textsuperscript{(12)} और
\VUgras{तेल-सिरके का दान} \textsuperscript{(13)} है, जिसमें सलाद सजाने का सब कुछ
रहता है : तेल, सिरका, \VUgras{नमक} \textsuperscript{(14)} और \VUgras{मिर्च}
\textsuperscript{(15)}। अंत में सबसे नीचे के ख़ाने में एक \VUgras{केक} \textsuperscript{(17)} है,
\VUgras{सरसों का दान} \textsuperscript{(16)} के पास।

%%K t06-c3-04-1 p
4. --- भोजन-कक्ष की घड़ी, जिसे \VUgras{दीवार-घड़ी} \textsuperscript{(20)} कहते
हैं, बहुत अनोखे आकार की है। वह लकड़ी के एक चौखटे में जड़ी है, जिसमें एक
\VUgras{वायुदाब-मापी} \textsuperscript{(19)} और एक \VUgras{तापमापी}
\textsuperscript{(18)} भी हैं।

%%K t06-tit-6 sub
\VUsaut{4.6mm}
\VUpk{10.2pt}{40}{भोजन परोसना।}
\VUsaut{3.0mm}

%%K t06-c3-05-1 p
5. --- कमरे के चारों ओर दीवारों पर मोम की पॉलिश की बलूत की
\VUgras{लकड़ी की चौखट} \textsuperscript{(21)} लगी है। कपड़े का एक
\VUgras{दरवाज़े का परदा} \textsuperscript{(22)} बरामदे की ओर जाते दरवाज़े को अच्छी
तरह ढक देता है। भोजन के बाद परिवार वहीं कॉफ़ी पीने जाएगा।
\VUgras{प्याले} \textsuperscript{(24)} और \VUgras{तश्तरियाँ} \textsuperscript{(25)}
पहले से \VUgras{चीनी बरतन की अलमारी} \textsuperscript{(23)} पर सजे हैं।

%%K t06-c3-06-1 p
6. --- मेज़ के ऊपर छत से एक \VUgras{लटकता लैंप} \textsuperscript{(26)} बँधा है;
उसकी बहुत तेज़ रोशनी शाम को सारे भोजन-कक्ष में भर जाती है।

%%K t06-c3-07-1 p
7. --- अब ख़ुद \VUgras{भोजन की मेज़} \textsuperscript{(27)} देखें। परिवार के आने
पर मेज़ लगी हुई थी। \VUgras{मेज़पोश} \textsuperscript{(28)} पर हर खाने वाले की जगह
रखी थी एक \VUgras{प्लेट} \textsuperscript{(33)}, एक \VUgras{नैपकिन}
\textsuperscript{(29)}, एक \VUgras{चम्मच} \textsuperscript{(34)}, एक \VUgras{काँटा}
\textsuperscript{(35)}, एक \VUgras{छुरी} \textsuperscript{(36)} और एक
\VUgras{गिलास} \textsuperscript{(32)}। शराब की \VUgras{बोतलें}
\textsuperscript{(30)} और पानी की एक \VUgras{सुराही} \textsuperscript{(31)} भी थी।

%%K t06-c3-08-1 p
8. --- \VUgras{नौकर} \textsuperscript{(39)} पहले \VUgras{सूप का पात्र}
\textsuperscript{(37)} लाया, जिसमें अब भी कुछ सूप से भाप उठ रही है। अब वह पहला
\VUgras{व्यंजन} \textsuperscript{(38)} परोस रहा है, माँस का। फिर वह वे व्यंजन
लाएगा जिनके नाम हम ले चुके; और अंत में, \VUgras{मिठाई} के बाद, वह इस बात का लाभ
उठाएगा कि मालिक बरामदे में चले गए हैं, और मेज़ ख़ाली करके भोजन-कक्ष फिर ठीक कर
देगा।

%%K t06-c3-08-2 p
\textit{टिप्पणी।} --- \textit{भोजन से जुड़े रोज़ के शब्द यहाँ केवल मोटे तौर पर
दिए गए हैं। यह सूची « होटल » (सं. 13) और « बाज़ार » (सं. 15) की दो तालिकाओं में
पूरी की जाएगी।}

%%K t06-tit-7 sub
\VUsaut{4.6mm}
\VUpk{10.2pt}{40}{रसोई।}
\VUsaut{3.0mm}

%%K t06-c4-01-1 p
1. --- अधखुले दरवाज़े से जिस रसोई पर हम नज़र डालते हैं, वहाँ एक सुरम्य
गड़बड़झाला मिलता है। \VUgras{चूल्हे} \textsuperscript{(1)} के सामने, जहाँ
\VUgras{आग} \textsuperscript{(3)} चटक रही है, \VUgras{भूनने का यंत्र}
\textsuperscript{(5)} लगा है। एक बढ़िया \VUgras{भुना माँस} \textsuperscript{(7)}
\VUgras{सींक} \textsuperscript{(6)} पर है और पकने को है।

%%K t06-c4-02-1 p
2. --- \VUgras{धौंकनी} \textsuperscript{(8)} और \VUgras{कोयले का बेलचा}
\textsuperscript{(9)}, जो अभी काम में आए हैं, \VUgras{लकड़ियों}
\textsuperscript{(10)} समेत ज़मीन पर ही छोड़ दिए गए हैं।

%%K t06-c4-03-1 p
3. --- अँगीठी के ऊपर का हिस्सा सँवरा है : \VUgras{लालटेन} \textsuperscript{(11)},
\VUgras{दियासलाई-दान} \textsuperscript{(13)} जिसमें \VUgras{दियासलाइयाँ}
\textsuperscript{(12)} हैं, \VUgras{शमादान} \textsuperscript{(14)} और
\VUgras{हाथ का शमादान} \textsuperscript{(15)} अपनी \VUgras{मोमबत्ती}
\textsuperscript{(16)} समेत — सब अपनी जगह हैं। पर बड़ी \VUgras{कोयले की टोकरी}
\textsuperscript{(18)}, जो उस \VUgras{कोयले} \textsuperscript{(17)} से भरी है जिसे
\VUgras{रसोइया} \textsuperscript{(23)} अभी तहख़ाने से लाई है, रसोई के बीच में ही
पड़ी रह गई है।

%%K t06-c4-04-1 p
4. --- सच तो यह है कि बेचारी को जल्दी करनी है, यदि दोपहर का भोजन समय पर तैयार
करना है। इस समय वह \VUgras{पकाने के चूल्हे} \textsuperscript{(19)} पर एक
\VUgras{रसा} \textsuperscript{(24)} चला रही है, जो जलने न पाए।

%%K t06-c4-05-1 p
5. --- \VUgras{तंदूर} \textsuperscript{(20)} में एक केक पक रहा है, जो उसने बच्चों
की शाम की चाय के लिए बनाया है। चूल्हे के ऊपर \VUgras{करछुल} \textsuperscript{(21)}
और \VUgras{झरना} \textsuperscript{(22)} टँगे हैं।

%%K t06-tit-8 sub
\VUsaut{4.0mm}
\VUpk{10.2pt}{40}{रसोई के बरतन।}
\VUsaut{3.0mm}

%%K t06-c4-06-1 p
6. --- कमरे के पीछे टँगा \VUgras{बरतनों का सेट} \textsuperscript{(25)} बहुत साफ़
है। ताँबे की सुंदर \VUgras{पतीलियाँ} \textsuperscript{(26)}, \VUgras{दूध का जग}
\textsuperscript{(27)}, \VUgras{छलनी} \textsuperscript{(28)} और \VUgras{कद्दूकस}
\textsuperscript{(29)} ध्यान से माँजे गए हैं।

%%K t06-c4-07-1 p
7. --- \VUgras{नमकदान} \textsuperscript{(30)} छोटी अलमारी पर \VUgras{चायदान}
\textsuperscript{(31)} और \VUgras{तेल की बड़ी बोतल} \textsuperscript{(32)} के पास
रखा है। \VUgras{दराज़} \textsuperscript{(33)} में शायद चम्मच-काँटे और रसोई की
छुरियाँ रहती हैं।

%%K t06-c4-08-1 p
8. --- \VUgras{झाड़ू} \textsuperscript{(34)} और \VUgras{पंख का झाड़न}
\textsuperscript{(35)} एक कोने में हैं। रसोइया अभी \VUgras{काटने का कुंदा}
\textsuperscript{(36)} काम में ला चुकी है; उसने उसे खिड़की के पास छोड़ दिया है।

%%K t06-c4-09-1 p
9. --- एक \VUgras{हाथ का तौलिया} \textsuperscript{(37)} दीवार पर टँगा है,
\VUgras{कीप} \textsuperscript{(38)} के पास। रसोई के इस भाग में \VUgras{सीख़-जाली}
\textsuperscript{(39)}, \VUgras{कटार} \textsuperscript{(40)} और
\VUgras{काटने का पटरा} \textsuperscript{(41)} रखे जाते हैं। फिर कटार के ऊपर, एक
\VUgras{ताक} \textsuperscript{(42)} पर, \VUgras{हाँडियाँ} \textsuperscript{(43)},
\VUgras{सालन की देगची} \textsuperscript{(44)} अपने \VUgras{ढक्कन}
\textsuperscript{(45)} समेत, और \VUgras{कड़ाहा} \textsuperscript{(46)} हैं।
\VUgras{तलने की कड़ाही} \textsuperscript{(47)} पास ही टँगी है।

%%K t06-c4-10-1 p
10. --- इस कोने में केवल ताँबे की \VUgras{टोंटी} \textsuperscript{(48)} चमक डालती
है। \VUgras{नाली-चौकी} \textsuperscript{(49)}, जिस पर बड़ा \VUgras{घड़ा}
\textsuperscript{(50)} रखा है, अपनी छोटी अलमारी के साथ बहुत काम की होगी, जहाँ
\VUgras{सिरके का पीपा} \textsuperscript{(51)} अपनी \VUgras{टोंटी}
\textsuperscript{(52)} समेत, \VUgras{कूड़े का संदूक} \textsuperscript{(53)} और
\VUgras{गंदे बरतन} \textsuperscript{(54)} रखे जाते हैं।

%%K t06-tit-9 sub
\VUsaut{4.0mm}
\VUpk{10.2pt}{40}{रसोई में रखी और चीज़ें।}
\VUsaut{3.0mm}

%%K t06-c4-11-1 p
11. --- वह बड़ा \VUgras{टब} \textsuperscript{(55)} जिसमें \VUgras{सब्ज़ियाँ}
\textsuperscript{(58)} धोई जाती हैं, और ढलवाँ लोहे का \VUgras{कड़ाहा} \textsuperscript{(56)},
जो \VUgras{ईंट के फ़र्श} \textsuperscript{(57)} पर रखा है, शायद उसी अलमारी के
हैं।

%%K t06-c4-12-1 p
12. --- रसोइया के पास चूल्हों की कमी नहीं, क्योंकि यह रहा मेज़ के कोने पर एक
\VUgras{गैस का चूल्हा} \textsuperscript{(59)}, जिस पर एक छोटी पतीली में पानी गरम
हो रहा है। उससे निकलती \VUgras{भाप} \textsuperscript{(60)} बताती है कि वह उबल रहा
होगा। यह पानी शायद कॉफ़ी के लिए है, क्योंकि मेज़ के दूसरे सिरे पर यह रहे
\VUgras{कॉफ़ी का बरतन} \textsuperscript{(64)} और \VUgras{कॉफ़ी की चक्की}
\textsuperscript{(63)}।

%%K t06-c4-13-1 p
13. --- गैस का चूल्हा \VUgras{गैस की टोंटी} \textsuperscript{(62)} से जुड़ा
है, एक लंबी \VUgras{रबर की नली} \textsuperscript{(61)} के ज़रिए।

%%K t06-c4-14-1 p
14. --- \VUgras{दिहाड़ी की मददगार} \textsuperscript{(66)}, जो रसोइया की सहायता को
आती है, रात के भोजन की सब्ज़ियाँ तैयार कर रही है। छीलकर और धोकर वह उन्हें
\VUgras{बरतन} \textsuperscript{(65)} में रख देगी, जब तक पकाने का समय न आए।

%%K t06-tit-10 sub
\VUsaut{4.0mm}
{\centering\fontsize{10.2pt}{10.2pt}\selectfont\textls[40]{\scshape स्नानघर।} \textit{(नहाने का कमरा।)}\par}
\VUsaut{3.0mm}

%%K t06-c5-01-1 p
1. --- घर के मुख्य कमरे जान लेने के लिए अब एक ही ताक-झाँक बाकी है : स्नानघर का
सामान देखना। इसमें देर न लगेगी, क्योंकि वह बड़ा नहीं, और हम झट देख लेंगे कि
\VUgras{नहाने के टब} \textsuperscript{(1)} के साथ एक अच्छा \VUgras{पानी गरम करने
का यंत्र} \textsuperscript{(2)} है, कि \VUgras{साबुनदानी} \textsuperscript{(3)}
नहाने वाले के हाथ की पहुँच में है, और कि \VUgras{तौलिया-डंडा}
\textsuperscript{(5)} से \VUgras{तौलिये} \textsuperscript{(4)} सुखाना सहज होता
होगा।

%%K t06-c5-02-1 p
2. --- इस कमरे की दीवारें \VUgras{चमकीली टाइलों} \textsuperscript{(6)} से ढकी
हैं, जिससे एक \VUgras{फ़ुहारा} \textsuperscript{(7)} लगाना संभव हुआ, बिना इस डर
के कि उसे बरतते समय छींटे उन्हें बिगाड़ दें। पैर धोने का जस्ते का एक छोटा
\VUgras{तसला} \textsuperscript{(8)} सामान पूरा कर देता है।
\VUsaut{6.0mm}
\VUfilet{22mm}
